% Copyright (c) 2026 Geovana Neves. Licensed under CC BY 4.0 (https://creativecommons.org/licenses/by/4.0/). See LICENSE-AND-AUTHORSHIP.md.
%
% overview/text/02-workflow-and-first-campaign.tex
%
% Parts 3 to 5 of guide 1: the workflow in one picture, a first campaign on
% the mock solver, and the map of the guides.

\section{The workflow}

\begin{frame}{Seven stages, one direction}
\begin{center}
\begin{tikzpicture}[
  st/.style={draw=hanblue,fill=hanblue!8,rounded corners=2pt,minimum height=9mm,
             text width=17mm,align=center,font=\scriptsize,inner sep=2pt},
  ar/.style={-{Stealth[length=5pt]},thick,hanblue}]
  \node[st] (a) at (0,0) {\textbf{1. Create}\\\code{init}, inputs};
  \node[st] (b) at (2.5,0) {\textbf{2. Plan}\\\code{plan}};
  \node[st] (c) at (5,0) {\textbf{3. Run}\\\code{run}};
  \node[st] (d) at (7.5,0) {\textbf{4. Follow}\\\code{status}, \code{log}};
  \node[st] (e) at (10,0) {\textbf{5. Collect}\\\code{collect}, \code{post}};
  \node[st,text width=30mm] (f) at (7.5,-1.6) {\textbf{6. Correct the records}\\\code{mark-failed}, \code{restore}};
  \node[st,text width=30mm] (g) at (10.9,-1.6) {\textbf{7. Manage storage}\\\code{free-space}, \code{sync}};
  \draw[ar] (a) -- (b); \draw[ar] (b) -- (c); \draw[ar] (c) -- (d); \draw[ar] (d) -- (e);
  \draw[ar,SoftGray,dashed] (d.south) -- (f.north);
  \draw[ar,SoftGray,dashed] (e.south) -- (e.south|-f.north);
\end{tikzpicture}
\end{center}
\begin{itemize}\small
  \item Stages 1 to 5 are every campaign. Stages 6 and 7 are for the day a
        record is wrong or the disk is full.
  \item \code{cfd-workspace} does stage 1. \code{cfd-matrix} does the rest.
  \item The cheatsheet, guide 4, has one page per stage.
\end{itemize}
\end{frame}

\begin{frame}[fragile]{The run matrix: one row per simulation}
\begin{lstlisting}
POL  | RUN | FLIGHT_CONDITION                  | SWEEP_VALUES | GEOMETRY | REF | SET | PPROC
-----|-----|-----------------------------------|--------------|----------|-----|-----|------
5001 | 1   | MACH:0.85, REmi:5, ALPHA:sweep    | 2.5,2.75     | wb.su2   | 001 | 001 | 001
5002 | 0   | MACH:0.85, REmi:5, ALPHA:3.0      | -            | wb.su2   | 001 | 002 | 001
\end{lstlisting}
\begin{itemize}\small
  \item \code{inputs/matrices/<name>.fs}: text, columns separated by
        \code{|}. A cell holding \code{-} states nothing.
  \item \code{POL} names the simulation. \code{RUN} 1 runs the row, 0 keeps
        it in the file without running it.
  \item \code{FLIGHT_CONDITION} holds \code{KEY:value} pairs. One key may say
        \code{sweep}; its values are in \code{SWEEP_VALUES}, one
        \textbf{point} each.
  \item \code{SET}, \code{REF} and \code{PPROC} name the three input files
        of the row.
\end{itemize}
\end{frame}

\begin{frame}{Three input files per row: s, r, p}
\relax{\cfdbbodysize
\begin{tabularx}{\linewidth}{@{}P{0.12\linewidth}P{0.38\linewidth}R@{}}
\toprule
\textbf{Column} & \textbf{File} & \textbf{It holds} \\
\midrule
\code{SET} & \code{inputs/setups/s<id>.toml} & how the solver runs: model,
  numerics, iteration limit. The keys belong to the solver. \\
\code{REF} & \code{inputs/references/r<id>.toml} & reference area, reference
  length and moment origin. \\
\code{PPROC} & \code{inputs/pproc/p<id>.toml} & what to keep from each run
  beyond the coefficients. \\
\code{GEOMETRY} & \code{inputs/mesh/<file>} & the mesh or geometry, in the
  format the solver reads. \\
\bottomrule
\end{tabularx}}
\vspace{2mm}
\begin{itemize}\small
  \item \code{001} and \code{s001} in the \code{SET} cell both name
        \code{s001.toml}.
  \item Many rows may share one file. A value the row and a file both state
        must agree.
\end{itemize}
\begin{takeaway}You write only under \code{inputs/}. Nothing a command writes
goes there.\end{takeaway}
\end{frame}

\begin{frame}{Plan first, then run}
\begin{columns}[T]
\begin{column}{0.48\textwidth}
\begin{block}{\code{cfd-matrix plan <matrix>}}
\begin{itemize}\small
  \item Reads the matrix and every file it names.
  \item Lists each point and the files its run would write.
  \item Launches nothing.
  \item Saves the plan receipt, \code{post/<matrix>/plan.json}.
\end{itemize}
\end{block}
\end{column}
\begin{column}{0.48\textwidth}
\begin{block}{\code{cfd-matrix run <matrix>}}
\begin{itemize}\small
  \item Needs the receipt.
  \item Refuses when an input changed after the plan: plan again.
  \item Writes one folder per point under \code{sims/} and one record per
        point in \code{runs.json}.
\end{itemize}
\end{block}
\end{column}
\end{columns}
\begin{takeaway}A mistake in an input is found by \code{plan}, in seconds,
and not by a job that waited in a queue.\end{takeaway}
\end{frame}

\begin{frame}{Records, and nothing deleted}
\begin{itemize}
  \item \code{runs.json} holds one record per point: its inputs and their
        hashes, its status, its times and the verdict on its convergence.
  \item A point that is already recorded \textbf{refuses} a second run.
        \code{--resume} runs only what is not recorded yet, and says how to
        redo a failed point; \code{--force-rerun} redoes a point you name.
  \item Before a record or a result is replaced, the old one moves to
        \code{archive/<stamp>/}.
  \item The commands that correct records or manage storage only
        \textbf{preview} until you add \code{--apply}.
\end{itemize}
\begin{takeaway}Finished evidence is never replaced by accident.\end{takeaway}
\end{frame}

\begin{frame}{Three verdicts}
\relax{\cfdbbodysize
\begin{tabularx}{\linewidth}{@{}P{0.22\linewidth}R@{}}
\toprule
\textbf{Verdict} & \textbf{Meaning} \\
\midrule
\code{converged} & the solver ended and its own evidence of convergence is
  there \\
\code{iteration_limit} & the solver ended properly without that evidence;
  the numbers exist and are not settled \\
\code{failed} & the process failed, or its output is missing, inconsistent
  or not a number \\
\bottomrule
\end{tabularx}}
\vspace{2mm}
\begin{itemize}\small
  \item The verdict is beside every row of the collected table.
  \item You may overrule it, with a reason on record:
        \code{cfd-matrix mark-failed} and \code{cfd-matrix mark-converged}.
\end{itemize}
\end{frame}

\section{A first campaign}

\begin{frame}[fragile]{Four small files}
\begin{lstlisting}
# inputs/matrices/first.fs
POL  | RUN | FLIGHT_CONDITION      | SWEEP_VALUES | REF | SET | PPROC
1001 | 1   | MACH:0.2, ALPHA:sweep | 0,2,4        | 001 | 001 | 001
\end{lstlisting}
\begin{columns}[T]
\begin{column}{0.36\textwidth}
\begin{lstlisting}[style=cfdbtoml]
# inputs/setups/s001.toml
[mock]
outcome = "converged"
iterations = 10
\end{lstlisting}
\begin{lstlisting}[style=cfdbtoml]
# inputs/pproc/p001.toml
history = true
\end{lstlisting}
\end{column}
\begin{column}{0.6\textwidth}
\begin{lstlisting}[style=cfdbtoml]
# inputs/references/r001.toml
area = 1.5
length = 1.0
moment_origin = [0, 0, 0]
axis_system = "body"
axis_directions = ["aft", "right", "up"]
\end{lstlisting}
\end{column}
\end{columns}
\vspace{1mm}
{\small One simulation, \code{1001}, with three points: angle of attack 0, 2
and 4 degrees at Mach 0.2.\par}
\end{frame}

\begin{frame}[fragile]{Six commands}
\begin{lstlisting}
cfd-workspace init campaign --solver mock
cd campaign
   (write the four files of the previous slide)
cfd-matrix plan first
cfd-matrix run first
cfd-matrix status --points
cfd-matrix collect first
\end{lstlisting}
\begin{itemize}\small
  \item \code{plan} lists three points. \code{run} runs them on this machine.
  \item \code{status --points} shows each point with its status and verdict.
  \item \code{collect} writes \code{post/first/campaign_sweep.csv}: one row
        per point, with the flight condition, the air state, the verdict and
        the coefficients.
\end{itemize}
\begin{takeaway}Change \code{outcome} to \code{"failed"} in a second setup
file and run a second row: you have seen a failure before it matters.\end{takeaway}
\end{frame}

\begin{frame}{From the mock to your solver}
\begin{enumerate}
  \item Make a new workspace: \code{cfd-workspace init <folder> --solver
        <name>}.
  \item Put the mesh in \code{inputs/mesh/} and name it in the
        \code{GEOMETRY} cell.
  \item Write the setup file with the solver's keys: guide 3 has one worked
        setup per solver.
  \item Say where the solver is on this machine: in the setup's
        \code{[solver]} table, or in \code{inputs/executables.toml} with the
        \code{SOLVER_BUILD} column. A run refuses before it records anything
        when the solver is missing.
  \item Plan, run, follow, collect: the same commands as on the mock.
\end{enumerate}
\end{frame}

\section{The guides}

\begin{frame}{The four guides}
\relax{\cfdbbodysize
\begin{tabularx}{\linewidth}{@{}P{0.05\linewidth}P{0.37\linewidth}R@{}}
\toprule
 & \textbf{File} & \textbf{Read it when} \\
\midrule
1 & \code{cfdb-guide-01-overview.pdf} & first: the idea, the solvers, the
  workflow \\
2 & \code{cfdb-guide-02-workspaces.pdf} & you start a campaign: the tree,
  the inputs, every stage, a workstation and a cluster \\
3 & \code{cfdb-guide-03-solvers.pdf} & you write a setup file or read a
  result: one part per solver \\
4 & \code{cfdb-guide-04-cheatsheet.pdf} & you are at the keyboard: every
  command and option, one page per stage \\
\bottomrule
\end{tabularx}}
\vspace{2mm}
\begin{cfdbevidence}[FlightStream]
A FlightStream workspace is run by \pyfsname{} [4]. Its commands, inputs and
guides are those of \pyfsname: \url{\pyfssite}.
\end{cfdbevidence}
\end{frame}
